\begin{abstract}
Noisy Partial Label Learning (NPLL) requires learning from candidate label sets that may contain false positives and, in the noisy case, may even exclude the ground-truth class. Many PLL and NPLL methods obtain stronger supervision by purifying, correcting, reconstructing, or reweighting candidate labels, but model-induced decisions that are carried across epochs can become persistent components of the working supervision. We propose \textbf{NSE} (\textbf{N}on-Destructive \textbf{S}upervision \textbf{E}xtraction), a framework based on \emph{epoch-wise source restoration}. NSE separates the native weak-supervision record from the epoch-local working supervision: at the beginning of each epoch, the working prior is restored from the original candidate mask or crowd prior, and all denoising, disambiguation, topology propagation, pseudo-label extraction, reliable selection, and salvage are applied only to temporary evidence. NSE uses Topology-DAES, a topology-guided Dynamic Adaptive Entropy Similarity kernel, for supervision extraction. Topology-DAES combines masked-entropy topology reliability with entropy-adaptive KNN propagation, allowing consistent neighborhoods to propagate sharper evidence while suppressing ambiguous weak-label neighborhoods. Model predictions are introduced only as temporary semantic evidence: they are gradually warmed up, projected by the native prior, and fused with the Topology-DAES belief through sample-wise reliability arbitration. Reliable selection and consensus-based salvage update training eligibility, not the persistent native record. Across synthetic NPLL, CIFAR-100H, and real crowdsourced datasets, NSE achieves strong performance without cross-epoch source write-back; reset-versus-write-back stress tests further show that persistent source insertion causes large source drift and wrong-write rates. These results suggest that NPLL robustness can be obtained by extracting reliable epoch-local supervision from noisy weak labels, rather than by carrying model-induced source modifications forward.
\end{abstract}

